\documentclass[12pt,a4paper]{article}
\usepackage[utf8]{inputenc}
\usepackage[T1]{fontenc}
\usepackage[spanish,es-nodecimaldot]{babel}
\usepackage{amsmath,amssymb,bm}
\usepackage{siunitx}
\usepackage{geometry}
\usepackage{booktabs}
\usepackage{array}
\usepackage{listings}
\usepackage{float}
\usepackage{tikz}
\usetikzlibrary{arrows.meta,positioning}
\geometry{margin=2.5cm}
\title{Modelo y estrategia de control de tres VSI en paralelo\\con MPC local y control secundario distribuido por consenso}
\author{}
\date{}
\lstset{basicstyle=\ttfamily\small,breaklines=true,frame=single,columns=fullflexible,showstringspaces=false}
\begin{document}
\maketitle

\section{Arquitectura general}
La microred está formada por tres inversores fuente de tensión (VSI) conectados en paralelo mediante filtros LCL a un punto común de acoplamiento (PCC). Cada inversor dispone de un controlador MPC local y el control secundario se incorpora mediante un algoritmo distribuido basado en consenso. Esta estructura coincide con la arquitectura planteada en el documento base, donde las variaciones de la carga y de los demás inversores se consideran perturbaciones externas desde el punto de vista del controlador local \cite{tesis}.

\section{Modelo eléctrico local}
Para cada inversor:
\begin{equation}
x_i=\begin{bmatrix}i_{f_i}&v_{C_i}&i_{c_i}\end{bmatrix}^{T},\qquad i=1,2,3.
\end{equation}
El modelo continuo es
\begin{equation}
\dot{x}_i=Ax_i+Bu_i+Ev_{PCC},
\end{equation}
con
\begin{equation}
A=\begin{bmatrix}0&-\frac{1}{L_f}&0\\\frac{1}{C_f}&0&-\frac{1}{C_f}\\0&\frac{1}{L_c}&0\end{bmatrix},\quad
B=\begin{bmatrix}\frac{1}{L_f}\\0\\0\end{bmatrix},\quad
E=\begin{bmatrix}0\\0\\-\frac{1}{L_c}\end{bmatrix}.
\end{equation}
Los parámetros empleados son
\begin{center}
\begin{tabular}{lc}\toprule
$L_f$&$1\,\mathrm{mH}$\\
$L_c$&$0.5\,\mathrm{mH}$\\
$C_f$&$5\,\mu\mathrm{F}$\\
$R$&$1\,\Omega$\\
$L$&$5\,\mathrm{mH}$\\
$V_{dc}$&$400\,\mathrm{V}$\\
$T_s$&$50\,\mu\mathrm{s}$\\
$f$&$60\,\mathrm{Hz}$\\
\bottomrule\end{tabular}
\end{center}

\section{Modelo del PCC y perturbaciones locales}
El voltaje del PCC es
\begin{equation}
v_{PCC}=\frac{R(i_{c1}+i_{c2}+i_{c3})+\frac{L}{L_c}(v_{C1}+v_{C2}+v_{C3})}{1+3\frac{L}{L_c}}.
\end{equation}
Esta expresión demuestra que el acoplamiento ocurre mediante $v_{C_i}$ e $i_{c_i}$ \cite{tesis}. Definiendo
\begin{equation}
\alpha=\frac{L}{L_c},\qquad K_{PCC}=\frac{1}{1+3\alpha},
\end{equation}
y
\begin{equation}
C_{PCC}=K_{PCC}\begin{bmatrix}0&\alpha&R\end{bmatrix},
\end{equation}
el modelo local utilizado por cada MPC es
\begin{equation}
\dot{x}_i=(A+EC_{PCC})x_i+Bu_i+Ed_i,
\end{equation}
\begin{equation}
v_{PCC}=C_{PCC}x_i+d_i.
\end{equation}
Las perturbaciones son
\begin{align}
d_1&=K_{PCC}[\alpha(v_{C2}+v_{C3})+R(i_{c2}+i_{c3})],\\
d_2&=K_{PCC}[\alpha(v_{C1}+v_{C3})+R(i_{c1}+i_{c3})],\\
d_3&=K_{PCC}[\alpha(v_{C1}+v_{C2})+R(i_{c1}+i_{c2})].
\end{align}

\section{Discretización y MPC}
Con $T_s=50\,\mu s$:
\begin{equation}
x_i[k+1]=A_dx_i[k]+B_du_i[k]+E_dd_i[k],
\end{equation}
\begin{equation}
y_i[k]=C_dx_i[k]+D_du_i[k].
\end{equation}
El MPC utiliza
\begin{equation}
N_p=20,\qquad N_u=5.
\end{equation}
La variable manipulada es
\begin{equation}
u_i=v_{inv,i}^{*},
\end{equation}
con
\begin{equation}
-200\,\mathrm{V}\leq u_i\leq200\,\mathrm{V}.
\end{equation}
Para una referencia de $120\,V_{rms}$:
\begin{equation}
V_{pico}=\sqrt{2}(120)\approx169.7\,V.
\end{equation}
El documento base establece la discretización a $50\,\mu s$ y la estructura de horizonte móvil del MPC \cite{tesis}.

\section{PWM y VSI}
La salida del MPC es una tensión promedio:
\begin{equation}
MPC_i\rightarrow u_i=v_{inv,i}^{*}\rightarrow PWM\rightarrow VSI_i.
\end{equation}
Para modulación bipolar de dos niveles:
\begin{equation}
m_i=\frac{2u_i}{V_{dc}},\qquad -1\leq m_i\leq1,
\end{equation}
\begin{equation}
D_i=\frac{1+m_i}{2}=\frac{1}{2}+\frac{u_i}{V_{dc}}.
\end{equation}

\section{Control secundario distribuido por consenso}
El objetivo nominal es
\begin{equation}
V^{*}=120\,V_{rms}.
\end{equation}
Cada inversor posee una corrección secundaria $\delta_i$:
\begin{equation}
\boxed{V_i^{ref}=V^{*}+\delta_i}.
\end{equation}
La referencia instantánea para cada MPC es
\begin{equation}
\boxed{r_i(t)=\sqrt{2}V_i^{ref}\sin(2\pi60t)}.
\end{equation}

\subsection{Medición local}
A partir del segundo estado:
\begin{equation}
v_{C_i}=x_i(2),
\end{equation}
se calcula
\begin{equation}
\boxed{V_{i,local}=\mathrm{RMS}(v_{C_i})}.
\end{equation}
Para $60\,Hz$, una ventana de aproximadamente un ciclo es
\begin{equation}
T_{RMS}=\frac{1}{60}=16.667\,ms.
\end{equation}
El error local es
\begin{equation}
e_i=V^{*}-V_{i,local}.
\end{equation}

\section{Topología de comunicación}
Para hacer visible el carácter distribuido se emplea inicialmente:
\begin{equation}
VSI_1\longleftrightarrow VSI_2\longleftrightarrow VSI_3.
\end{equation}
La matriz de adyacencia es
\begin{equation}
\boxed{
A_g=\begin{bmatrix}0&1&0\\1&0&1\\0&1&0\end{bmatrix}}.
\end{equation}
La matriz de grados es
\begin{equation}
D_g=\begin{bmatrix}1&0&0\\0&2&0\\0&0&1\end{bmatrix},
\end{equation}
y la Laplaciana:
\begin{equation}
\boxed{
L=D_g-A_g=
\begin{bmatrix}1&-1&0\\-1&2&-1\\0&-1&1\end{bmatrix}}.
\end{equation}

\section{Ley de consenso}
Definiendo
\begin{equation}
\delta=\begin{bmatrix}\delta_1&\delta_2&\delta_3\end{bmatrix}^{T},\qquad
e=\begin{bmatrix}e_1&e_2&e_3\end{bmatrix}^{T},
\end{equation}
la ley implementada es
\begin{equation}
\boxed{\dot{\delta}=k_Ie-k_cL\delta}.
\end{equation}
Para la topología en cadena:
\begin{align}
\dot{\delta}_1&=k_Ie_1+k_c(\delta_2-\delta_1),\\
\dot{\delta}_2&=k_Ie_2+k_c[(\delta_1-\delta_2)+(\delta_3-\delta_2)],\\
\dot{\delta}_3&=k_Ie_3+k_c(\delta_2-\delta_3).
\end{align}
Con Euler:
\begin{equation}
\boxed{\delta[k+1]=\delta[k]+T_s(k_Ie[k]-k_cL\delta[k])}.
\end{equation}
Se emplean inicialmente
\begin{equation}
k_I=5,\qquad k_c=10,
\end{equation}
y
\begin{equation}
-10\,V\leq\delta_i\leq10\,V.
\end{equation}
Las condiciones iniciales son
\begin{equation}
\delta_1(0)=\delta_2(0)=\delta_3(0)=0.
\end{equation}

\section{Interpretación de $\delta_i$}
El consenso busca
\begin{equation}
\delta_1\rightarrow\delta^{*},\qquad
\delta_2\rightarrow\delta^{*},\qquad
\delta_3\rightarrow\delta^{*}.
\end{equation}
No es necesario que $\delta^{*}=0$. En condiciones nominales, si las mediciones locales son $120\,V_{rms}$ y las correcciones parten de cero, entonces $\delta_i=0$ y $V_i^{ref}=120\,V_{rms}$. Ante una perturbación, $\delta_i$ puede apartarse de cero y luego converger entre sí.

\section{Generación de la senoide en Simulink}
Se utilizan bloques estándar: \texttt{Sine Wave}, \texttt{Gain} con ganancia $\sqrt{2}$ y \texttt{Product}. El bloque Sine Wave genera
\begin{equation}
s(t)=\sin(2\pi60t),
\end{equation}
y el producto entrega
\begin{equation}
r_i(t)=(\sqrt{2}V_i^{ref})s(t).
\end{equation}
Para $V_i^{ref}=120\,V$:
\begin{equation}
r_i(t)=169.7\sin(2\pi60t),
\end{equation}
correspondiente idealmente a $120\,V_{rms}$.

\section{Código MATLAB del consenso}
\begin{lstlisting}[language=Matlab]
function [V1ref,V2ref,V3ref,delta1,delta2,delta3] = ...
    consenso(V1local,V2local,V3local)

Ts = 50e-6;
Vnom = 120;

kI = 5;
kc = 10;

delta_min = -10;
delta_max = 10;

persistent delta_1 delta_2 delta_3

if isempty(delta_1)
    delta_1 = 0;
    delta_2 = 0;
    delta_3 = 0;
end

% Error local: e_i = Vnom - V_i_local
e_1 = Vnom - V1local;
e_2 = Vnom - V2local;
e_3 = Vnom - V3local;

% Topologia:
% VSI 1 <------> VSI 2 <------> VSI 3

Ag = [0 1 0;
      1 0 1;
      0 1 0];

Dg = [1 0 0;
      0 2 0;
      0 0 1];

% Laplaciana: L = Dg - Ag
L = Dg - Ag;

delta = [delta_1;
         delta_2;
         delta_3];

e = [e_1;
     e_2;
     e_3];

% Ley: delta_dot = kI*e - kc*L*delta
delta_dot = kI*e - kc*L*delta;

% Euler
delta = delta + Ts*delta_dot;

% Saturacion
delta(1) = min(max(delta(1),delta_min),delta_max);
delta(2) = min(max(delta(2),delta_min),delta_max);
delta(3) = min(max(delta(3),delta_min),delta_max);

delta_1 = delta(1);
delta_2 = delta(2);
delta_3 = delta(3);

% Referencias
V1ref = Vnom + delta_1;
V2ref = Vnom + delta_2;
V3ref = Vnom + delta_3;

delta1 = delta_1;
delta2 = delta_2;
delta3 = delta_3;

end
\end{lstlisting}

\section{Estructura final}
\begin{equation}
\boxed{
v_{C_i}\rightarrow RMS\rightarrow V_{i,local}\rightarrow Consenso
\rightarrow V_i^{ref}\rightarrow Senoidal\ 60Hz\rightarrow MPC_i
\rightarrow u_i\rightarrow PWM\rightarrow VSI_i\rightarrow LCL_i\rightarrow PCC
}
\end{equation}

El objetivo final se verifica mediante
\begin{equation}
\boxed{V_{PCC,rms}\approx120\,V_{rms}}.
\end{equation}
El consenso no sustituye al MPC: modifica las referencias de amplitud y el MPC continúa realizando el control primario rápido. Para evidenciar el transitorio del consenso se recomienda aplicar una perturbación, como un cambio de carga, y observar simultáneamente $V_{PCC,rms}$, $V_{i,local}$, $\delta_i$ y $V_i^{ref}$.

\section{Observación sobre la garantía}
La generación de la referencia garantiza que $V_i^{ref}=120\,V_{rms}$ produzca una senoide de amplitud pico $169.7\,V$. No constituye, por sí sola, una garantía matemática de que el PCC alcance exactamente $120\,V_{rms}$ para cualquier perturbación. La validación debe realizarse observando la respuesta de $V_{PCC,rms}$ y demostrando convergencia para las condiciones de operación estudiadas.

\begin{thebibliography}{9}
\bibitem{tesis}
Documento base del proyecto, \textit{MPC\_Microgrid.pdf}. Modelo de tres VSI en paralelo, filtro LCL, acoplamiento mediante PCC, discretización a $T_s=50\,\mu s$ y diseño del MPC.
\end{thebibliography}

\end{document}
